\documentclass[11pt]{article}
\usepackage{amsmath,amssymb,amsthm}
\newtheorem{theorem}{Theorem}
\begin{document}
\title{An Exact Feature Bound for Separable Reciprocal Queries}
\author{}
\date{November 24, 2026}
\maketitle

\paragraph{Scope and verification.}
This is a complete written sign-pattern bound for a fixed-query row
family. It is not Lean checked. In particular, it is not an unrestricted
HeadComplexity lower bound. Its two-head case is sharp, by the existing
five-query construction in \texttt{reverse-joint-five-row.tex}, and also
by the arbitrarily compact five-query construction. Its general bound is
exponential in the number of heads, so it does not exclude a
superconstant saving in the requested local compiler.

Fix real query data $D_{hj},b_{hj}$, with $1\leq h\leq H$ and
$1\leq j\leq q$. The variable row parameters are $a_h,p_h$, restricted
only by $D_{hj}-p_h>0$ for every $h,j$. Consider the strict signs of
\[
 F_j(a,p)=\sum_{h=1}^H\frac{a_h+b_{hj}}{D_{hj}-p_h}.
 \tag{1}
\]
There is no separately variable row bias in (1). A fixed query-dependent
bias would change the feature calculation, as would adding a separately
variable row bias. Neither operation is silently included below.

\begin{theorem}
For $H\geq1$, the family (1) has strict VC dimension at most
$2^{H+1}-3$. In particular, one head cannot shatter two queries and two
heads cannot shatter six queries.
\end{theorem}
\begin{proof}
For $S\subseteq[H]$, put $p_S=\prod_{h\in S}p_h$, with $p_\varnothing=1$.
For nonempty $S$, put
\[
 A_S=\sum_{h\in S}a_h p_{S\setminus\{h\}}.
\]
Multiplying (1) by its positive denominator product gives
\[
 P_j=\sum_h(a_h+b_{hj})\prod_{k\ne h}(D_{kj}-p_k).
\]
Expanding the products and grouping the terms containing an $a_h$ gives
\[
 \begin{split}
 P_j={}&\sum_{\varnothing\ne S\subseteq[H]}
       (-1)^{|S|-1} A_S\prod_{k\notin S}D_{kj}\\
 &+\sum_{S\subsetneq[H]}(-1)^{|S|}p_S
       \sum_{h\notin S}b_{hj}
          \prod_{k\notin S\cup\{h\}}D_{kj}.
 \end{split}                                                    \tag{2}
\]
Indeed, an $a_h p_T$ term occurs precisely when $h\notin T$;
setting $S=T\cup\{h\}$ gives the first line. The remaining terms give
the second line directly.

Consequently every $P_j$ is an affine function of the common features
\[
 (A_S)_{\varnothing\ne S\subseteq[H]},\qquad
 (p_S)_{\varnothing\ne S\subsetneq[H]}.
\]
There are $(2^H-1)+(2^H-2)=d:=2^{H+1}-3$ features.
An arbitrary feature image, including the one here, can visit no more
strict sign patterns of $q$ affine functions on $\mathbb R^d$ than the
whole space does.

For completeness, $q$ affine hyperplanes in $\mathbb R^d$ have at most
$\sum_{i=0}^d\binom qi$ complementary regions. Adding the final
hyperplane splits only regions met by it, and its intersections with
previous hyperplanes form an arrangement in dimension at most $d-1$.
Induction therefore gives the recurrence
$R(q,d)\leq R(q-1,d)+R(q-1,d-1)$ and the stated binomial bound.
Every nonempty strict sign cell is an intersection of open half-spaces,
hence convex and connected, so there is at most one cell for each
region and no extra sign patterns. Degenerate or constant affine
functions only decrease this count. If $q=d+1$, the bound is
$2^q-1$, strictly less than all $2^q$ sign patterns. Shattering a larger
set would shatter a subset of size $d+1$, which is impossible.
\end{proof}

\paragraph{The sharp two-head calculation.}
For $H=2$, the five features can be written
\[
 a_1,\quad a_2,\quad p_1,\quad p_2,\quad
 z=a_1p_2+a_2p_1.
\]
The exact cleared query is
\[
 P_j=-z+D_{2j}a_1+D_{1j}a_2-b_{2j}p_1-b_{1j}p_2
       +b_{1j}D_{2j}+b_{2j}D_{1j}.
 \tag{3}
\]
Thus six queries give six affine hyperplanes in five-dimensional feature
space, with at most $63$ strict sign cells. The parameter image lies on
the graph $z=a_1p_2+a_2p_1$, but no geometric property of that graph is
needed for this upper bound. The five-query lower construction therefore
attains the largest possible number of shattered queries in (1).
The higher-head bounds given by the same argument are $13$ for three
heads and $29$ for four heads. They are not asserted to be sharp.

\paragraph{What remains open.}
The feature count does not establish an $O(H)$ VC bound. The degree and
number of features grow with $H$, and the positive-denominator parameter
region can have many algebraic fibers without acquiring correspondingly
many orthants. Controlling the strict sign patterns of this family as
$H$ grows is still open in this run. Even an $O(H)$ result here would
need a separate argument before implying a lower bound for all original
HeadComplexity representations.

The companion exact checker evaluates (2) at rational parameters for
$1\leq H\leq7$, checks positivity of all denominators, and checks the
special two-head formula (3). The proof of (2) itself is the product
expansion above, not an inference from those tests.
\end{document}
